\chapter{Problem}\label{ch:problem}

extensive explanations of the problem at hand and why this needs to be improved 

\section{Repetitive Steps in the Iterative Workflow}
\begin{figure}[h]
    \centering
    \includegraphics[width=\linewidth]{FlowChart_IterativeProcess.png}
    \caption{Iterative workflow in 3D printing highlighting the step to be optimized}
    \label{workflow}
\end{figure}

\section{Negative Effects on Professionals and (casual) Makers}

``If the user’s print failed or did not match their design intent, they would have to adjust either their model or print settings to correct the issue, often with the operator’s assistance''\cite[p.~389]{hudson_understanding_2016}

\section{what does it solve}

\begin{itemize}
    \item what does the optimized workflow look like? 
    \item how it might make sharing files easier for beginners when supports and modifiers are included $\rightarrow$ lowering the access barrier for casuals, beginners and in educational contexts
    \item how it reduces cognitive workload $\rightarrow$ minimizing mistakes and waste due to prints going wrong 
\end{itemize}

Definition of supports 

``With FDM printers, geometric constraints relating to overhangs and base size can be partially addressed using automatically generated scaffolding known as supports and rafts''\cite[p.~391]{hudson_understanding_2016}

``many print centers enforced final check-overs by the operators on users’ designs and print settings before printing to reduce failed prints''\cite[p.~392]{hudson_understanding_2016}